% Method section — Feature Engineering subsection (Part C)
% Paste this into reports/midsem/sections/method.tex, alongside
% Part A's label-construction subsection (method_label_construction.tex).
%
% Numbers here come from the pooled two-sequence easy feature tables (MH_01_easy,
% MH_02_easy) produced by scripts/make_features.py and the analysis in notebooks/04.
% If the feature set or sequences change, re-run that notebook and update them.

\subsection{Feature Engineering}
\label{sec:feature-engineering}

Each window is a $200 \times 6$ matrix of raw IMU samples (three gyroscope
and three accelerometer axes over one second). Classic regressors cannot
consume this raw tensor directly, so we summarise every window as a single
fixed-length vector of $104$ hand-crafted features. The guiding hypothesis
is that open-loop inertial integration accumulates position error fastest in
windows that are \emph{energetic}, \emph{erratic}, or \emph{abrupt}; the
feature families below are designed to quantify exactly those properties from
the signal alone, using only raw IMU (no ground-truth information leaks in).

\subsubsection{Time-domain features (60)}
For each of the six channels we compute ten statistics: mean, standard
deviation, variance, RMS, minimum, maximum, peak-to-peak range, skewness,
excess kurtosis, and zero-crossing rate. The first-moment and energy terms
(mean, RMS, variance) measure how much the drone moved and how large a bias
the integrator must contend with; the shape terms (skewness, kurtosis) flag
asymmetric or spike-laden windows; and the zero-crossing rate, computed about
each channel's own mean so the accelerometer gravity offset does not dominate,
is a cheap proxy for oscillation frequency. This family contains the single
strongest linear correlate of the label, the accelerometer $x$-axis mean and
RMS ($|r| \approx 0.67$), and supplies the top-ranked features under a
RandomForest importance analysis (\texttt{ax\_mean} and \texttt{ax\_rms}).

\subsubsection{Frequency-domain features (18)}
Time-domain statistics describe \emph{how much} a window moves but not
\emph{how} it moves. For each channel we remove the per-window mean, take a
real FFT, and extract the dominant frequency, the total spectral energy, and
the normalised spectral entropy of the power spectrum. Mean removal is
essential: without it the accelerometer's large gravity/DC term would be read
as a spurious $0\,\mathrm{Hz}$ peak. Spectral entropy distinguishes clean,
tonal motion (energy in one band) from broadband, noise-like windows; the
per-axis spectral-entropy features are among the strongest label correlates
($|r| \approx 0.47$--$0.52$), consistent with noise-like windows drifting
more than smoothly manoeuvring ones.

\subsubsection{Cross-axis and magnitude features (26)}
A drone is a rigid body, so its axes are coupled; three groups of features
capture coupling the per-axis statistics miss. Accelerometer and gyroscope
\emph{vector magnitudes} (mean, standard deviation, maximum of
$\lVert a \rVert$ and $\lVert \omega \rVert$) give an orientation-independent
measure of motion, invariant to the drone's heading. \emph{Jerk} features
(per-axis RMS and the mean and peak magnitude of $\mathrm{d}a/\mathrm{d}t$)
target abrupt changes, which are precisely where discrete strapdown
integration between samples is least accurate. Finally, the fifteen pairwise
Pearson \emph{correlations} between channels describe coordinated motion:
strong cross-axis correlation is the signature of a banking turn or climb
rather than uncorrelated sensor noise, and several accelerometer--gyroscope
cross terms (e.g.\ \texttt{corr\_gz\_ax}) rank highly under RandomForest
importance.

\subsubsection{Scaling and the no-leakage constraint}
Because the project splits by sequence (leave-one-sequence-out,
Section~\ref{sec:label-construction}), feature standardisation must respect
fold boundaries: a \texttt{StandardScaler} is fit on each fold's training
sequences only and then applied to the held-out sequence. Fitting a scaler
across all data first would leak test-set statistics into training. The
handoff tables \texttt{features\_<seq>.csv} are therefore stored
\emph{unscaled}, and scaling is performed per fold at modelling time.

\subsubsection{Redundancy and feature selection}
The feature set is deliberately broad and contains known collinearities
(variance is the square of standard deviation; RMS approximates the absolute
mean when the mean dominates). Pooled across the two easy sequences, $22$ feature
pairs exceed a $0.95$ absolute correlation. Tree-based models are insensitive
to this redundancy, and a feature-selection experiment
(notebook~\texttt{04\_feature\_engineering}) confirms that a compact subset of
the highest-ranked features retains most of the predictive signal; the full
set is retained for Phase~1 so that per-family contributions can be compared
honestly in the ablation discussion.
